%% Lean: NeoTiling.Planar.{Chain,FlowLP,Farkas}.
\section{Dimension two: solvability is a linear program}\label{sec:planar}

Throughout this section $d=2$, $P=\set{p_t\in\Rpp}_{t\in[T]}$, $X_t:=p_t^1$ and $Y_t:=p_t^2$. Strict solvability (\cref{thm:solvability}) and weak solvability (\cref{thm:closure}) are treated in parallel.

\begin{definition}[mode $\star$]\label{def:modes}
In the \emph{strict mode} $u\gs0$ means $u>0$ and $\beta_\star:=0$; in the \emph{weak mode} $u\gs0$ means $u\ge0$ and $\beta_\star:=1$. For $P=\set{p_t\in\Rp[d]}_{t\in[T]}$ with any $d\ge2$, a set $S\subset[T]$ belongs to $\Spm(P)$ if for every $t\in S$ there is a nonzero $\xi\in\Rp[d]$ with $\ip{\xi}{p_s-p_t}\gs0$ for all $s\notin S$. Thus $\Spm(P)=\Sp(P)$ (\cref{def:reachable}) in the strict mode and $\Spm(P)=\wSp(P)$ (\cref{def:weak}) in the weak mode. The problem $\WS_\star(d)$: given rational $P\subset\Rpp[d]$ and $y\in\Q^T_+$, decide whether $y\in\cone\set{\ind S\mid S\in\Spm(P)}$.
\end{definition}

\subsection{Chains}

For $a,b,q\in[T]$ put $\chi(a,b,q):=(X_b-X_a)(Y_q-Y_a)-(Y_b-Y_a)(X_q-X_a)$. If $X_a<X_b$, then $\chi(a,b,q)>0$ if and only if $p_q$ lies strictly above the line through $p_a$ and $p_b$.

\begin{definition}[chain]\label{def:chain}
A \emph{chain} is a sequence $V=(v_0,\dots,v_k)$ of indices with $X_{v_0}<\dots<X_{v_k}$, $Y_{v_0}>\dots>Y_{v_k}$ and $\chi(v_j,v_{j+1},v_{j+2})>0$ for $j+2\le k$. The \emph{region} $C(V)$ is the set of $q\in[T]$ with
\[
X_q\ge X_{v_0},\qquad Y_q\ge Y_{v_k},\qquad \chi(v_j,v_{j+1},q)\ge0\quad(j<k),
\]
and the \emph{interior} $C^\circ(V)$ is the set of $q$ satisfying all these inequalities strictly. The \emph{forbidden region} is $F_\star(V):=C(V)$ in the strict mode and $F_\star(V):=C^\circ(V)$ in the weak mode.
\end{definition}

\begin{lemma}[hull chain]\label{lem:hull}
Let $\varnothing\ne U\subset[T]$. Let $v_0$ be the lexicographic minimum of $(X,Y)$ on $U$ and, while some $q\in U$ has $Y_q<Y_{v_i}$, let $v_{i+1}$ be the rightmost among these points of least slope as seen from $v_i$. Then $V=(v_0,\dots,v_k)$ is a chain with vertices in $U$ and $U\subset C(V)$.
\end{lemma}

\begin{proof}
A point $q\in U$ with $Y_q<Y_{v_i}$ has $X_q>X_{v_i}$: for $i=0$ by the choice of $v_0$; for $i\ge1$, by induction $X_q>X_{v_{i-1}}$, and $X_q\le X_{v_i}$ would give $q$ a smaller slope than $v_i$ as seen from $v_{i-1}$. Hence $X$ increases and $Y$ decreases along $V$, and the construction stops. The slope from $v_j$ to $v_{j+2}$ is at least the slope to $v_{j+1}$, and equality would contradict the choice of the rightmost point; so $\chi(v_j,v_{j+1},v_{j+2})>0$. Let $q\in U$. Then $X_q\ge X_{v_0}$ and, as the construction stopped, $Y_q\ge Y_{v_k}$. For $\chi(v_j,v_{j+1},q)\ge0$: if $Y_q<Y_{v_j}$, then $q$ was a candidate at step $j$ and its slope from $v_j$ is not smaller; if $Y_q\ge Y_{v_j}$ and $X_q\ge X_{v_j}$, then $q$ lies on or above the line through $v_j,v_{j+1}$, which descends to the right; if $X_q<X_{v_j}$, then $j\ge1$, $q$ lies on or above the line through $v_{j-1},v_j$ by induction on $j$, and to the left of $v_j$ this line lies above the line through $v_j,v_{j+1}$, whose slope is larger.
\end{proof}

\begin{theorem}[chains]\label{thm:chain}
A set $S\subset[T]$ belongs to $\Spm(P)$ if and only if $S=[T]$ or there is a chain $V$ such that no vertex of $V$ lies in $S$, every index outside $C(V)$ lies in $S$, and no index of $F_\star(V)$ lies in $S$.
\end{theorem}

\begin{proof}
\emph{Sufficiency.} For $S=[T]$ take $\xi=(1,1)$. Otherwise let $t\in S$. All $s\notin S$ lie in $C(V)$, and since $t\notin F_\star(V)$, one of the inequalities defining $C(V)$ fails at $t$ in the sense $\gs$: $X_{v_0}-X_t\gs0$, or $Y_{v_k}-Y_t\gs0$, or $-\chi(v_j,v_{j+1},t)\gs0$. Take $\xi=(1,0)$, $\xi=(0,1)$ or the edge normal $\xi=(Y_{v_j}-Y_{v_{j+1}},X_{v_{j+1}}-X_{v_j})\ge0$, respectively. Then $\ip{\xi}{p_s-p_t}$ equals $(X_s-X_{v_0})+(X_{v_0}-X_t)$, $(Y_s-Y_{v_k})+(Y_{v_k}-Y_t)$ or $\chi(v_j,v_{j+1},s)-\chi(v_j,v_{j+1},t)$; the first summand is nonnegative as $s\in C(V)$, and the second is $\gs0$.

\emph{Necessity.} Let $S\ne[T]$ and $U:=[T]\setminus S$. By \cref{lem:hull} there is a chain $V$ with vertices in $U$ and $U\subset C(V)$. Suppose $q\in F_\star(V)\cap S$ with witness $\xi$; then $\ip{\xi}{p_{v_i}-p_q}\gs0$ for all $i$. If $X_q>X_{v_k}$, then $p_q\ge p_{v_k}$, with strict inequality in both coordinates in the weak mode, so $\ip{\xi}{p_{v_k}-p_q}\ngs0$. If $X_q\le X_{v_0}$, the mode is strict (as $q\in C^\circ(V)$ forces $X_q>X_{v_0}$), $X_q=X_{v_0}$ and $Y_q\ge Y_{v_0}$, so $\ip{\xi}{p_{v_0}-p_q}\le0$. Otherwise $X_{v_j}<X_q\le X_{v_{j+1}}$ for some $j$; with $\lambda:=X_{v_{j+1}}-X_q\ge0$, $\mu:=X_q-X_{v_j}>0$, $A:=\ip{\xi}{p_{v_j}-p_q}$ and $B:=\ip{\xi}{p_{v_{j+1}}-p_q}$,
\[
\lambda A+\mu B=-\xi_2\,\chi(v_j,v_{j+1},q)\le0 .
\]
In the strict mode this contradicts $A,B>0$. In the weak mode the sum vanishes while $\chi(v_j,v_{j+1},q)>0$, so $\xi_2=0$, $\xi_1>0$ and $A=\xi_1(X_{v_j}-X_q)<0$, a contradiction.
\end{proof}

\subsection{The chain graph}

\begin{definition}[chain graph]\label{def:graph}
The graph $G_P$ has the nodes $\sigma$ (source), $\tau$ (sink), $A_b$ for $b\in[T]$ (``the chain starts at $b$'') and $E_{ab}$ for $a,b\in[T]$ with $X_a<X_b$, $Y_a>Y_b$ (``the last edge is $ab$''). The arcs are
\[
\sigma\to\tau,\quad \sigma\to A_b,\quad A_a\to E_{ab},\quad E_{ab}\to E_{bc}\ \text{ if }X_b<X_c,\ Y_b>Y_c,\ \chi(a,b,c)>0,\quad A_b\to\tau,\quad E_{ab}\to\tau.
\]
Every arc $e$ carries two sets of indices $\Out_e$ and $\Bnd_e$.
\begin{itemize}
\item $\sigma\to A_b$ (left slab of $b$): $\Out_e=\set{q\ne b\mid X_q<X_b\ \text{or}\ (X_q=X_b,\ Y_q<Y_b)}$ and $\Bnd_e=\set{q\ne b\mid X_q=X_b,\ Y_q\ge Y_b}$.
\item $A_a\to E_{ab}$ and $E_{ca}\to E_{ab}$ (slab of the edge $ab$): $\Out_e=\set{q\ne b\mid X_a<X_q\le X_b,\ \chi(a,b,q)<0}$ and $\Bnd_e=\set{q\ne b\mid X_a<X_q\le X_b,\ \chi(a,b,q)=0}$.
\item $A_b\to\tau$ and $E_{ab}\to\tau$ (right slab of $b$): $\Out_e=\set{q\mid X_q>X_b,\ Y_q<Y_b}$ and $\Bnd_e=\set{q\mid X_q>X_b,\ Y_q=Y_b}$.
\item $\sigma\to\tau$: $\Out_e=[T]$ and $\Bnd_e=\varnothing$.
\end{itemize}
For a path $\pi$ from $\sigma$ to $\tau$ put $\Out(\pi):=\bigcup_{e\in\pi}\Out_e$ and $\Bnd(\pi):=\bigcup_{e\in\pi}\Bnd_e$.
\end{definition}

The graph is acyclic: $\mathrm{rk}(\sigma):=0$, $\mathrm{rk}(A_b):=1$, $\mathrm{rk}(E_{ab}):=2+\abs{\set{q\mid X_q<X_b}}$, $\mathrm{rk}(\tau):=T+3$ strictly increases along arcs.

\begin{theorem}[family]\label{thm:family}
$S\in\Spm(P)$ if and only if $\Out(\pi)\subseteq S\subseteq\Out(\pi)\cup\beta_\star\Bnd(\pi)$ for some path $\pi$ from $\sigma$ to $\tau$ in $G_P$, where $\beta_\star\Bnd(\pi)$ is $\Bnd(\pi)$ in the weak mode and $\varnothing$ in the strict mode. Moreover, for every such path the sets $\Out_e\cup\Bnd_e$, $e\in\pi$, are pairwise disjoint and $\Out(\pi)\cap\Bnd(\pi)=\varnothing$, so that
\begin{equation}\label{eq:disjoint}
\ind{\Out(\pi)}=\sum_{e\in\pi}\ind{\Out_e},\qquad \ind{\Bnd(\pi)}=\sum_{e\in\pi}\ind{\Bnd_e}.
\end{equation}
\end{theorem}

\begin{proof}
By the list of arcs, the paths from $\sigma$ to $\tau$ are $\sigma\to\tau$ and the paths $\sigma\to A_{v_0}\to E_{v_0v_1}\to\dots\to E_{v_{k-1}v_k}\to\tau$ of the chains $V=(v_0,\dots,v_k)$. The path $\sigma\to\tau$ has $\Out(\pi)=[T]$ and corresponds to $S=[T]$. For the path of a chain $V$, the slabs $X_q\le X_{v_0}$, $X_{v_j}<X_q\le X_{v_{j+1}}$ and $X_q>X_{v_k}$ of its arcs partition $[T]$, and the sets of each arc lie in its slab; this gives the disjointness and \eqref{eq:disjoint}. Inside a slab, membership in $C(V)$ is local. In the slab of the edge $j$ put $\lambda:=X_{v_{j+1}}-X_q$, $\mu:=X_q-X_{v_j}$ and $c:=\chi(v_j,v_{j+1},q)$; for every edge $i$,
\[
(\lambda+\mu)\,\chi(v_i,v_{i+1},q)=\lambda\,\chi(v_i,v_{i+1},v_j)+\mu\,\chi(v_i,v_{i+1},v_{j+1})+c\,(X_{v_{i+1}}-X_{v_i}),
\]
and likewise $(\lambda+\mu)(Y_q-Y_{v_k})=\lambda(Y_{v_j}-Y_{v_k})+\mu(Y_{v_{j+1}}-Y_{v_k})+c$, while $X_q>X_{v_j}\ge X_{v_0}$. The first two terms are nonnegative because the slopes along a chain are monotone and $Y$ decreases along it, so $q\in C(V)$ if and only if $c\ge0$, and $q\in C^\circ(V)$ if and only if $c>0$. In the left slab $q\in C(V)$ if and only if $X_q=X_{v_0}$ and $Y_q\ge Y_{v_0}$, and $q\notin C^\circ(V)$; in the right slab $q\in C(V)$ (resp.\ $C^\circ(V)$) if and only if $Y_q\ge Y_{v_k}$ (resp.\ $>$). Comparing with \cref{def:graph},
\[
\Out(\pi)=[T]\setminus C(V),\qquad \Bnd(\pi)=\bigl(C(V)\setminus C^\circ(V)\bigr)\setminus V,
\]
and the vertices of $V$ lie in $C(V)\setminus C^\circ(V)$. Hence the three conditions of \cref{thm:chain} are equivalent to the two inclusions: in the strict mode they say $S=\Out(\pi)$, in the weak mode that $S$ contains $\Out(\pi)$ and no vertex or interior point.
\end{proof}

\subsection{The flow LP}

\begin{lemma}[cube]\label{lem:cube}
Let $O\cap B=\varnothing$ and $0\le z\le\ind B$. Then $\ind O+z\in\cone\set{\ind{O\cup B'}\mid B'\subset B}$.
\end{lemma}

\begin{proof}
The cube $[0,1]^B$ is the convex hull of the indicators of subsets of $B$; equivalently, by induction on $\abs B$, $\ind O+z=(1-z_q)(\ind O+z|_{B_0})+z_q(\ind{O\cup\set q}+z|_{B_0})$ for $B=B_0\cup\set q$.
\end{proof}

\begin{theorem}[$d=2$]\label{thm:lp}
A vector $y\in\R^T$ lies in $\cone\set{\ind S\mid S\in\Spm(P)}$ if and only if the system
\[
f\ \text{is a $\sigma$--$\tau$ flow in }G_P,\qquad y=\sum_ef_e\ind{\Out_e}+g,\qquad 0\le g\le\beta_\star\sum_ef_e\ind{\Bnd_e}
\]
is feasible.
\end{theorem}

\begin{proof}
$(\Rightarrow)$ Let $y=\sum_Sm_S\ind S$ with $m_S\ge0$, $S\in\Spm(P)$. By \cref{thm:family} choose paths $\pi_S$ with $\Out(\pi_S)\subseteq S\subseteq\Out(\pi_S)\cup\beta_\star\Bnd(\pi_S)$ and put $f:=\sum_Sm_S\ind{\pi_S}$, $g:=\sum_Sm_S\ind{S\setminus\Out(\pi_S)}$. Then $f$ is a flow, and \eqref{eq:disjoint} gives $y=\sum_ef_e\ind{\Out_e}+g$ and $0\le g\le\beta_\star\sum_ef_e\ind{\Bnd_e}$.

$(\Leftarrow)$ Since $G_P$ is acyclic, $f=\sum_im_i\ind{\pi_i}$ with positive weights $m_i>0$ and $\sigma$--$\tau$ paths $\pi_i$ \citep[Thm.~3.5]{AhujaMagnantiOrlin1993}. By \eqref{eq:disjoint}, $\sum_ef_e\ind{\Out_e}=\sum_im_i\ind{\Out(\pi_i)}$ and $h:=\sum_ef_e\ind{\Bnd_e}=\sum_im_i\ind{\Bnd(\pi_i)}$. Split $g$ proportionally: $z_i(q):=g_q/h_q$ for $q\in\Bnd(\pi_i)$ (then $h_q\ge m_i>0$) and $z_i(q):=0$ otherwise. Then $0\le z_i\le\beta_\star\ind{\Bnd(\pi_i)}$ and $y=\sum_im_i(\ind{\Out(\pi_i)}+z_i)$. By \cref{lem:cube} each summand lies in the cone of the vectors $\ind{\Out(\pi_i)\cup B'}$ with $B'\subset\beta_\star\Bnd(\pi_i)$, and these sets belong to $\Spm(P)$ by \cref{thm:family}.
\end{proof}

\begin{corollary}[polynomiality]\label{cor:poly}
In both modes $\WS_\star(2)$ is solvable in polynomial time, and for $y$ in the cone a representation $y=\sum_km_k\ind{S_k}$ with $S_k\in\Spm(P)$ and at most $T$ terms is found in polynomial time.
\end{corollary}

\begin{proof}
The graph $G_P$ is built in $O(T^4)$ arithmetic operations on numbers of degree two in the input. The system of \cref{thm:lp} has $O(T^3)$ variables, $O(T^2)$ constraints, coefficients $0,\pm1$ and right-hand side $y$, so its feasibility is decided in polynomial time \citep[Ch.~13--15]{Schrijver1986}. A solution yields a representation via flow decomposition and \cref{lem:cube} as in the proof of \cref{thm:lp}, and the number of terms is reduced to $T$ by Carath\'eodory's theorem for conic hulls \citep[Cor.~17.1.2]{Rockafellar1970}.
\end{proof}

\subsection{The column problem and the certificate}

For $c\in\R^T$ and $S\subset[T]$ put $c(S):=\sum_{q\in S}c_q$. The \emph{weight} of an arc $e$ of $G_P$ is $c(\Out_e)+\beta_\star\sum_{q\in\Bnd_e}\min(c_q,0)$, and the weight $\mathrm{wt}_c(\pi)$ of a path is the sum over its arcs.

\begin{corollary}[column problem]\label{cor:col}
Let $c\in\R^T$. Then $\min_{S\in\Spm(P)}c(S)$ equals the minimum weight of a $\sigma$--$\tau$ path in the acyclic graph $G_P$, and $c\ge0$ on $\Spm(P)$ if and only if every such path has nonnegative weight.
\end{corollary}

\begin{proof}
By \eqref{eq:disjoint}, $\mathrm{wt}_c(\pi)=c(\Out(\pi))+\beta_\star\sum_{q\in\Bnd(\pi)}\min(c_q,0)$. This value is $c(S)$ for $S:=\Out(\pi)\cup\set{q\in\beta_\star\Bnd(\pi)\mid c_q<0}$, which lies in $\Spm(P)$ by \cref{thm:family}. Conversely, every $S\in\Spm(P)$ has the form $\Out(\pi)\cup B'$ with $B'\subseteq\beta_\star\Bnd(\pi)$ by \cref{thm:family}, and then $c(S)\ge\mathrm{wt}_c(\pi)$. Shortest paths in an acyclic graph are computed in linear time \citep[Ch.~4]{AhujaMagnantiOrlin1993}.
\end{proof}

\begin{corollary}[certificate of unsolvability]\label{cor:farkas}
$y\notin\cone\set{\ind S\mid S\in\Spm(P)}$ if and only if there is $c\in\R^T$ with $\ip cy<0$ and $c(S)\ge0$ for all $S\in\Spm(P)$; equivalently, with $\ip cy<0$ and nonnegative weight of every $\sigma$--$\tau$ path in $G_P$. The certificate is checked by one shortest-path computation.
\end{corollary}

\begin{proof}
The cone is finitely generated, hence polyhedral and closed \citep[Cor.~7.1a]{Schrijver1986}. The first equivalence is the Farkas lemma \citep[Cor.~7.1d]{Schrijver1986}; the dual condition over the family $\Spm(P)$ is a shortest-path condition by \cref{cor:col}.
\end{proof}
